\documentclass[
    ngerman,
    accentcolor=3b,
    % dark_mode, % uncomment for dark mode
    fontsize= 12pt,
    a4paper,
    aspectratio=169,
    colorback=true,
    fancy_row_colors,
    leqno,
    fleqn,
    boxarc=3pt,
    fleqn,
    main,
    % shell_escape = false, % Kompatibilität mit sharelatex
]{algoslides}

%%------------%%
%%--Packages--%%
%%------------%%

% \usepackage{audutils}
% \usepackage{fopbot}
\usepackage{booktabs}
\usepackage{etoolbox}


%%----------------------------%%
%%--Stilistische Anpassungen--%%
%%----------------------------%%


\renewcommand\tabularxcolumn[1]{m{#1}}% for vertical centering text in X column
% Remove unwanted space from tables
\aboverulesep = 0mm \belowrulesep = 0mm
\renewcommand{\arraystretch}{1.4}

%%---------------------------%%
%%--Dokumenteneinstellungen--%%
%%---------------------------%%

% Workshop-Individuelle Einstellungen
\def\workshoptitle{<Workshoptitel>}
\def\shortworkshoptitle{<Kurztitel>}
\subtitle{<Subtitel>}
\def\gruesswoerte{<Begrüßende Worte>}
\author{<Autor1>\and<Autor2>}
\titlegraphic*{\includegraphics{example-image}}


% Ophase-Einstellungen
\ExplSyntaxOn
\def\summerophasename{Sommerophase}
\def\winterophasename{Winterophase}
\DeclareDocumentCommand{\ophase}{}{
    \int_compare:nTF {\the\month > 4}{
        \winterophasename\space{}\the\year-\the\numexpr\year+1\relax
    }{
        \summerophasename\space{}\the\year
    }
}
\ExplSyntaxOff
\graphicspath{{pictures}}
\title[\shortworkshoptitle{}]{\workshoptitle{}}
\department{TU Darmstadt | Fachbereich Informatik | \ophase{} | \insertshorttitle{}}
\date{\today}
\logo*{\includegraphics{d120-logo.png}}

%%-------------------------%%
%%--Beginn des Dokumentes--%%
%%-------------------------%%

\begin{document}

    %%-----------%%
    %%--Titelei--%%
    %%-----------%%

    \maketitle{}

    \begin{frame}[c]
        % Die Begrüßenden Worte können individuell pro Workshop festgelegt werden. Am Besten nur was kurzes, sowie "Gude", "Hi", "Herzlich Willkommen!", ...
        \centering\huge\textbf{\gruesswoerte{}}
    \end{frame}

    %%---------------------------%%
    %%--Beginn der Präsentation--%%
    %%---------------------------%%

    \begin{frame}
        \frametitle{Das steht heute auf dem Plan}
        \tableofcontents[subsubsectionstyle=hide]
    \end{frame}


    %%-----------%%
    %%--Kapitel--%%
    %%-----------%%

    % !TeX root = chapter01.tex
\documentclass[
    ngerman,
    accentcolor=3b,
    % dark_mode, % uncomment for dark mode
    fontsize= 12pt,
    a4paper,
    aspectratio=169,
    colorback=true,
    fancy_row_colors,
    leqno,
    fleqn,
    boxarc=3pt,
    fleqn,
    main,
    design=2008,
    % shell_escape = false, % Kompatibilität mit sharelatex
]{algoslides}
\RequirePackage{import}
\subimport{../common}{preamble}

%%--------------------------%%
%%--Imports from Main File--%%
%%--------------------------%%

% Get Labels from Main Document using the xr-hyper Package
\externaldocument[ext:]{../main}
% Set Graphics Path, so pictures load correctly
\graphicspath{{../pictures}}

\begin{document}
    \section{Einführung}\label{1}\label{Einfuehrung}
    \subsection{Subsection1}
    \begin{frame}[<+(1)->]
        \slidehead{}
        \begin{itemize}
            \item Item 1
            \item Item 2
            \item Item 3
        \end{itemize}
    \end{frame}

    \subsection{Subsection2}
    \begin{frame}
        \slidehead{}
        More text
    \end{frame}
\end{document}

\end{document}
